\section{Conclusion}
\label{sec:conclusion}

This paper began with a Bi-level Adversarial Reinforcement Learning (Bi-ARL) formulation for network threat detection and then evolved into a broader bilevel adversarial training framework. The reproduced NSL-KDD results still support a narrow positive claim for the original RL formulation: Bi-ARL improves over vanilla PPO and offers the best mean F1 among the RL baselines. However, the more important modern-data finding is that the refactored route-C models are substantially stronger than the old RL line. In particular, BiAT-FTTransformer becomes the strongest neural detector on UNSW-NB15 and consistently improves over BiAT-MLP, although it still remains behind the best boosted-tree baselines. The current contribution should therefore be framed as a competitive bilevel adversarial training framework with a stronger tabular backbone, not as a state-of-the-art IDS. The most important next steps are to exploit multi-GPU parallelism for larger sweeps, further stabilize the inner loop, strengthen baseline coverage, and extend the evaluation to harder protocols and additional modern datasets such as CSE-CIC-IDS2018.
